\section{当前定位、限制与前瞻评价}
\label{sec:limits}
\subsection{已经建立的结论}
项目已经有可独立运行的离线 Python 静态审查交付路径：使用者明确选择源码根，系统保留字节身份、解析分母、候选、已登记事实、有限性质状态、缺项和一致产物。共享内核能够拒绝已知的错 subject、弱静态确认、晚到同范围冲突和嵌套别名漂移。保留研究系统可以组织有界检索、规划、类型化工具、有限 ReAct、单专家票、quorum 与条件执行见证。

在少量已选开发候选和给定分类场景中，见证仍定位为\textbf{有界确认工具}。其正向控制和 K2 迁移等价性已经建立；自主场景提出、合格外部迁移、公开入口以及 fresh held-out 泛化没有由当前证据建立。原 G-A=NO 与新产品 GP0/GP1 的工程通过能够同时成立，回答的是不同问题。

\begin{longtable}{L{33mm}L{70mm}L{46mm}}
\caption{当前边界、已观察影响和下一步证据要求。}\\
\toprule 边界 & 当前证据 & 需要补的证据\\\midrule\endfirsthead
\caption[]{（续）}\\
\toprule 边界 & 当前证据 & 需要补的证据\\\midrule\endhead
静态语义 & 未支持 helper、动态分派、外部库及对象协议保持 UNKNOWN；有限模型有正例 & 对明确语义逐项新增正反控制，不宣称整个 Python 完备。\\
性质与全应用 & 有界 REFUTED/SUPPORTED 明确不覆盖应用安全 & 另行定义资产、路径、输入域和安全性质。\\
源码与环境 & 产品身份覆盖所选 Python 清单；静态 execution digest 为 null & 条件执行需可信环境清单和单独准入。\\
场景选择 & E-iter 四格没有第二物理场景；两格因分支，两格因环境 & 独立冻结场景研究、反馈、重复次数和预算。\\
公开入口 & callback 直接调用没有公开 parser；本地调用不建立远程权限 & 显式入口/部署范围及授权依据。\\
文件系统写入 & 历史隔离阻止临时目录，sink 尚未到达 & T10 私有 scratch 的拥有者夹具和写入边界控制。\\
原生扩展 & 3b 暂定外部候选未完成对应环境准备 & T11 可信源码、工具链、hash 和可复核构建。\\
多专家与早停 & 可共享一份原观察；固定已知票保证有限 & 独立 live 增量与成本实验；未知未来证据仍是边界。\\
数据盲性 & 开发、oracle 和暴露元数据不具 held-out 资格 & 新鲜、事前冻结、标签隔离且按仓库计分。\\
用户价值 & GP3/GP4 尚未测量 & 配对人工审查与可选模型说明的净增量。\\
\bottomrule\end{longtable}

\subsection{工程数值能说明什么}
2684 个通过用例说明当前声明控制没有在完整套件中失败；它不枚举所有程序输入，也不是 detector benchmark。31 条必需验收证明静态交付契约，未覆盖三个条件能力。七个 CLI 产物用例证明不同运行状态下可交付，并显示约 269 毫秒的小用例启动墙钟；无法外推大仓库吞吐。

历史 447812 tokens 的投入以及四格 ABSTAIN 说明本轮自主任务没有达到门槛；不能在缺少反事实对照时将原因唯一归为模型、预算或协议。F3 与 F4 的区分使后续研究能针对场景/环境分别设计验证。0 合格外部样本的含义是当前准入管线没有完成迁移测量，而不是估计某总体成功率为零。

\subsection{已冻结但未执行的增量评估协议}
T14 前瞻方案以仓库为统计单位，纳入新鲜项目拥有或明确授权源码，按语言、性质、源码规模和初始上下文缺口分层；排除 oracle、开发专用、已暴露 held-out 和未授权材料。源码、性质契约、基线事实与准入上下文事前冻结，评价标签不进入 detector prompts、路径别名或源码选择。\citep{increment}

配对臂为原人工流程、离线产品报告和条件模型说明。前两者项目模型请求为零；模型说明仍关闭，未设新请求或 token 授权。说明只能引用合法事实与缺项，不能修改性质状态。审阅顺序在审阅者间平衡，保留排除和失败项，记录人工秒数、事实错误、缺项识别、行动性、全部模型请求/usage、unknown usage 与端到端墙钟。

\begin{pseudocode}{alg:future}{尚未执行的前瞻 paired-review 协议}
\Pinput{事前冻结且明确授权的新鲜仓库、独立评价记录、reviewer 与预算}
\Poutput{仓库级配对质量/时间差及完整成本；允许无收益结果}
\Pstep{核对数据角色、授权、源码和契约身份；泄漏或漂移立即停止。}
\Pstep{先冻结同一基线事实和上下文，再运行原人工流程与离线产品报告臂。}
\Pstep{平衡 reviewer 的臂顺序，记录时间、事实错误、正确缺项及行动价值。}
\Pstep{仅在单独授权后启用模型说明；不允许说明改写性质状态。}
\Pstep{以仓库为单位汇总配对差，保留未完成/未知项及所有成本。}
\Pstep{按事前分析规则判断 GP3/GP4；无净收益同样为有效结果。}
\end{pseudocode}

若之后设计区间分析，应事前确定以仓库为单元的配对 bootstrap 或其他适用方法、样本量和缺失处理。当前冻结协议只确立仓库级配对汇总与完整记账，没有事后假装已完成显著性实验。planner、Jev、多专家或 target execution 进入产品需要独立可复现净收益，不能因为组件已存在就自动启用。

\subsection{延后路线和发布状态}
T10/T11 的条件能力未启用，T14 只有协议。历史 K5--K8 没有启动。定位报告列出的四条路线仍分别是私有临时目录、原生扩展源码构建、场景提出研究和新评测集；这些路线的设计依赖继续保留，启用时需选择具体范围并重新预登记。\citep{positioning,architectureplan}

当前修复已验收但尚未提交，HEAD 没有因这些修复产生新提交；工作区也不是 clean。报告属于工作树快照，源码聚合 SHA-256 与输入回执哈希用于复核。历史汇总 JSON 已修正 K2、G-A 和统计来源，旧值只保留在历史快照对象。没有借重写历史或追加模型实验改善本报告结论。
